\section{Physical covariant variation and conditional string extensions}
\label{sec:physical-variational-string}

The auxiliary action of Sect.~\ref{sec:variational-multisymplectic} is an
exact and useful structure, but it is not an action for the spacetime and
its matter. This section supplies the physical variational comparison. We
first formulate Einstein gravity coupled to vortical dust and identify its
covariant presymplectic and boundary data. We then ask a deliberately
separate question: what survives if the exterior theory is enlarged to
Einstein--Maxwell--dilaton--axion gravity? The answer includes a rigidity
theorem, but not a string-theoretic explanation of the rotating-dust
interior.

Throughout this section the spacetime signature is $(-,+,+,+)$,
$G_{\rm N}$ is Newton's constant, and all orientations and normal directions
are fixed before boundary quantities are compared.

\subsection{Einstein gravity with vortical dust}

A pressureless fluid is not adequately represented by one velocity
potential when vorticity is present. Brown's material-coordinate action
keeps the particle current and three comoving labels as independent fields
~\cite{Brown1993CQG}. Let $\mathcal J^\mu$ be a future-directed vector density,
let $\alpha^A$ be material labels, and let $\vartheta$ and $\beta_A$ be
Clebsch variables, where $A=1,2,3$. Define
\begin{equation}
 |\mathcal J|_g:=
 \sqrt{-g_{\mu\nu}\mathcal J^\mu\mathcal J^\nu},
 \qquad
 n:=\frac{|\mathcal J|_g}{\sqrt{-g}},
 \qquad
 u^\mu:=\frac{\mathcal J^\mu}{|\mathcal J|_g}.
 \label{eq:physical-Brown-variables}
\end{equation}
For particles of constant rest mass $m>0$, the dust action is
\begin{equation}
 \mathcal S_{\rm B}[g,\mathcal J,\vartheta,\alpha,\beta]
 =\int_{\mathcal M}
 \left[-m|\mathcal J|_g
 +\mathcal J^\mu
 \left(\partial_\mu\vartheta
 +\beta_A\partial_\mu\alpha^A\right)\right]\,d^4x.
 \label{eq:physical-Brown-action}
\end{equation}
This is the $\rho(n)=mn$ specialization of the perfect-fluid action. Its
independent variations give
\begin{align}
\partial_\mu\mathcal J^\mu&=0,
&
\mathcal J^\mu\partial_\mu\alpha^A&=0,
 \nonumber\\
\partial_\mu(\beta_A\mathcal J^\mu)&=0,
&
m u_\mu+\partial_\mu\vartheta
+\beta_A\partial_\mu\alpha^A&=0.
 \label{eq:physical-Brown-transport}
\end{align}
Metric variation produces
\begin{equation}
 T^{\rm d}_{\mu\nu}=mn\,u_\mu u_\nu.
 \label{eq:physical-Brown-stress}
\end{equation}
Moreover, Eq.~\eqref{eq:physical-Brown-transport} implies
\begin{equation}
 m\,d u^\flat=-d\beta_A\wedge d\alpha^A,
 \label{eq:physical-Brown-vorticity}
\end{equation}
so the action admits genuine vortical flow. In particular, the stationary
rotating dust considered in this paper need not be forced into an
irrotational scalar ansatz.
The constant shift $\vartheta\mapsto\vartheta+c$ has Noether current
$\mathcal J^\mu$, while material-label transformations preserving
$d\vartheta+\beta_A d\alpha^A$ encode the circulation symmetries. These are
matter symmetries, not replacements for spacetime diffeomorphism charges.

On a spacetime region with smooth non-null boundary components, the
physical Dirichlet action is
\begin{align}
 \mathcal S_{\rm phys}
 =&\frac{1}{16\pi G_{\rm N}}
 \int_{\mathcal M}R\,\epsilon_g
 +\frac{1}{8\pi G_{\rm N}}
 \int_{\partial\mathcal M}\epsilon_n K\,\epsilon_h
 +\mathcal S_{\rm B},
 \label{eq:physical-EH-GHY-Brown-action}
\end{align}
where $\epsilon_n=g(n,n)=\pm1$, $h$ is the induced metric, and
$K=h^{ab}K_{ab}$ in the convention of
Eq.~\eqref{eq:extrinsic-curvature-convention}. Corner terms must be added
when boundary components meet non-smoothly. The second term in
Eq.~\eqref{eq:physical-EH-GHY-Brown-action} is the
Gibbons--Hawking--York term in these conventions
~\cite{GibbonsHawking1977PRD}.

For the bulk Lagrangian four-form $\boldsymbol L_{\rm phys}$, write its
first variation as
\begin{equation}
 \delta\boldsymbol L_{\rm phys}
 =\boldsymbol E(\Psi)\cdot\delta\Psi
 +d\boldsymbol\theta(\Psi;\delta\Psi),
 \qquad
 \Psi=(g,\mathcal J,\vartheta,\alpha,\beta).
 \label{eq:physical-covariant-first-variation}
\end{equation}
The Lee--Wald presymplectic current and the diffeomorphism Noether current
are, respectively,
\begin{align}
 \boldsymbol\omega(\Psi;\delta_1\Psi,\delta_2\Psi)
 &:={\delta_1}\boldsymbol\theta(\Psi;\delta_2\Psi)
 -{\delta_2}\boldsymbol\theta(\Psi;\delta_1\Psi),
 \label{eq:physical-Lee-Wald-current}\\
 \boldsymbol J_\xi
 &:=\boldsymbol\theta(\Psi;\mathcal L_\xi\Psi)
 -\iota_\xi\boldsymbol L_{\rm phys}.
 \label{eq:physical-diffeomorphism-current}
\end{align}
For solutions and linearized solutions,
$d\boldsymbol\omega=0$. Diffeomorphism covariance gives
$\boldsymbol J_\xi=\boldsymbol C_\xi+d\boldsymbol Q_\xi$, with the
constraint term $\boldsymbol C_\xi$ vanishing on shell. Subject to the
chosen boundary conditions and integrability of the right-hand side, the
Hamiltonian variation is
\begin{equation}
 \delta H_\xi
 =\int_{\partial C}
 \left(\delta\boldsymbol Q_\xi
 -\iota_\xi\boldsymbol\theta\right)
 \label{eq:physical-Hamiltonian-variation}
\end{equation}
on a Cauchy surface $C$~\cite{LeeWald1990JMP,IyerWald1994PRD}. Boundary
counterterms and corners shift the representative of
$\boldsymbol\theta$ and must be included before a charge is declared
finite or integrable. Equations~\eqref{eq:physical-Lee-Wald-current}--\eqref{eq:physical-Hamiltonian-variation}
refer to the coupled metric--dust theory. They cannot be replaced by the
weighted meridional current of
Eq.~\eqref{eq:ms-current-form}.
In particular, choosing $\xi$ to be the stationary or axial Killing field
defines physical energy or angular-momentum variations only after the
asymptotic normalization, boundary conditions, and integrability test have
been supplied.

The GHY term also makes the relation between variational and Darmois data
precise. For a non-null three-boundary define the gravitational momentum
density conjugate to $h^{ab}$ by
\begin{equation}
 \Pi_{ab}:=
 \frac{\epsilon_n\sqrt{|h|}}{16\pi G_{\rm N}}
 \left(K_{ab}-Kh_{ab}\right).
 \label{eq:physical-boundary-momentum}
\end{equation}

\begin{proposition}[Darmois data as equality of boundary momenta]
\label{prop:physical-Darmois-momentum}
Let two $C^2$ spacetime regions meet across a non-null hypersurface
$\Sigma$, and express both second fundamental forms using one common normal
orientation. Once the first fundamental forms have been identified,
\begin{equation}
 [\Pi_{ab}]=0
 \quad\Longleftrightarrow\quad
 [K_{ab}]=0.
 \label{eq:physical-momentum-Darmois-equivalence}
\end{equation}
Thus a shell-free stationary point of the joined Dirichlet action has the
complete Darmois data. Equality only of selected Killing-block derivatives
does not suffice.
\end{proposition}

\begin{proof}
With $h_{ab}$ common, $[\Pi_{ab}]=0$ is equivalent to
$[K_{ab}]-h_{ab}[K]=0$. Contracting with the inverse three-metric gives
$[K]-3[K]=-2[K]=0$. Hence $[K]=0$, and substitution gives
$[K_{ab}]=0$. The converse is immediate. With the outward normals used in
the two separate actions, one first reverses one normal to obtain the
common-orientation convention. A nonzero momentum jump is the variational
counterpart of an Israel surface stress
~\cite{Darmois1927,Israel1966NCB,MarsSenovilla1993CQG}.
\end{proof}

There is one further reduction caveat. Substituting a stationary,
axisymmetric ansatz directly into
Eq.~\eqref{eq:physical-EH-GHY-Brown-action} and only then varying is valid
only when the relevant principle of symmetric criticality holds. Its group
cohomology and isotropy hypotheses are nontrivial in general relativity
~\cite{FelsTorre2002CQG}. The safe logical route in the present paper is to
derive the covariant equations and boundary variation first and then impose
the two Killing symmetries. No auxiliary scalar action establishes this
missing equivalence.

\subsection{Einstein-frame EMDA and neutral axidilaton rigidity}

We now change the field content. In a standard four-dimensional
Einstein-frame normalization, the bosonic EMDA bulk action is
\begin{align}
 \mathcal S_{\rm EMDA}^{E}
 =\frac{1}{16\pi G_{\rm N}}\int_{\mathcal M}\epsilon_E
 \bigg[&R_E-2(\nabla\varphi)^2
 -\frac12e^{4\varphi}(\nabla\kappa)^2
 -e^{-2\varphi}F_{\mu\nu}F^{\mu\nu}
 \nonumber\\
 &-\kappa F_{\mu\nu}{}^\star F^{\mu\nu}\bigg]
 +\mathcal S_{\rm B}[g_E].
 \label{eq:physical-EMDA-action}
\end{align}
Here $\varphi$ is the dilaton, $\kappa$ is the axion, $F=dA$, and the
dust is neutral and minimally coupled only to $g_E$. This convention is
equivalent, after the signature translation, to the primary EMDA actions
used in Refs.~\cite{GarciaGaltsovKechkin1995PRL,GaltsovKechkin1994PRD}.
The gravitational boundary term has the form already displayed in
Eq.~\eqref{eq:physical-EH-GHY-Brown-action} with $g=g_E$.

When $F=0$, the scalar equations are
\begin{equation}
 \nabla_\mu\!\left(e^{4\varphi}\nabla^\mu\kappa\right)=0,
 \qquad
 \Box_E\varphi=\frac12e^{4\varphi}(\nabla\kappa)^2.
 \label{eq:physical-neutral-axidilaton-equations}
\end{equation}
There is no dust source in either equation because
$\mathcal S_{\rm B}[g_E]$ is independent of $\varphi$ and $\kappa$.
For invariant scalars in the circular metric
\begin{equation}
 ds_E^2=G_{AB}dx^A dx^B
 +e^\mu(dr^2+dz^2),
 \qquad
 \det G=-r^2,
 \label{eq:physical-EMDA-circular-metric}
\end{equation}
Eq.~\eqref{eq:physical-neutral-axidilaton-equations} becomes
\begin{equation}
 \partial_a\!\left(re^{4\varphi}\partial_a\kappa\right)=0,
 \qquad
 \partial_a(r\partial_a\varphi)
 =\frac r2e^{4\varphi}|\nabla_{rz}\kappa|^2,
 \qquad a\in\{r,z\}.
 \label{eq:physical-neutral-orbit-equations}
\end{equation}

\begin{theorem}[Neutral axidilaton no-hair theorem]
\label{thm:physical-neutral-axidilaton-rigidity}
Let a connected, horizonless, stationary and axisymmetric solution of
Eq.~\eqref{eq:physical-EMDA-action} satisfy $F=0$. Assume all of the
following.

\begin{enumerate}
\item The spacetime is globally regular and asymptotically flat with one
end, the axial action has a regular fixed axis, and the complete orbit
space is a half-plane with no excised inner boundary. It admits one global
circular Weyl chart $(r,z)$ in which
Eq.~\eqref{eq:physical-EMDA-circular-metric} holds with
$\det G=-r^2$, the axis is $r=0$, and the orbit measure in the reduced
scalar equations is $r\,dr\,dz$.

\item The neutral dust is minimally coupled to $g_E$. Across every internal
matter--vacuum interface, $\varphi$, $\kappa$, and their canonical normal
fluxes are continuous; in particular there is no scalar surface source.

\item The invariant scalars are $C^2$ away from interfaces, $C^1$ in the
global weak sense, and regular at the axis. The axion has a globally
single-valued real lift, so there is no axionic monodromy.

\item With $R=(r^2+z^2)^{1/2}$, constants $\varphi_\infty$ and
$\kappa_\infty$ exist such that
\begin{align}
 \varphi-\varphi_\infty&=O(R^{-1}),
 &
 |\nabla_{rz}\varphi|&=O(R^{-2}),
 \nonumber\\
 \kappa-\kappa_\infty&=O(R^{-1}),
 &
 |\nabla_{rz}\kappa|&=O(R^{-2}).
 \label{eq:physical-axidilaton-falloff}
\end{align}
The weighted scalar energies
\begin{equation}
 \int r|\nabla_{rz}\varphi|^2\,dr\,dz,
 \qquad
 \int re^{4\varphi}|\nabla_{rz}\kappa|^2\,dr\,dz
 \label{eq:physical-axidilaton-energies}
\end{equation}
are finite over the orbit space.
\end{enumerate}

Then
\begin{equation}
 \kappa\equiv\kappa_\infty,
 \qquad
 \varphi\equiv\varphi_\infty.
 \label{eq:physical-axidilaton-rigidity-conclusion}
\end{equation}
Consequently, appending a neutral, source-free axidilaton sector cannot
dress a globally regular galaxy satisfying these hypotheses.
\end{theorem}

\begin{proof}
Let $D_R$ be the intersection of the orbit half-plane with a Euclidean
half-disk of radius $R$, with small half-disks around any intermediate
corners used only during the exhaustion. Multiply the first equation in
Eq.~\eqref{eq:physical-neutral-orbit-equations} by
$\kappa-\kappa_\infty$ and integrate by parts. The interface terms cancel
because both the scalar and its canonical normal flux are continuous.
Regularity makes $\kappa$ and its first derivatives bounded at the axis,
where the factor $r$ makes the boundary term vanish. On the outer
semicircle, Eq.~\eqref{eq:physical-axidilaton-falloff} gives
\begin{equation}
 \int_{\partial D_R}
 r e^{4\varphi}(\kappa-\kappa_\infty)
 \partial_n\kappa\,ds=O(R^{-1}).
 \label{eq:physical-axion-boundary-estimate}
\end{equation}
There is no horizon or other inner boundary. Passing to the limit, using
Eq.~\eqref{eq:physical-axidilaton-energies}, yields
\begin{equation}
 \int_D re^{4\varphi}|\nabla_{rz}\kappa|^2\,dr\,dz=0.
 \label{eq:physical-axion-energy-zero}
\end{equation}
The integrand is nonnegative and $e^{4\varphi}>0$, so connectedness gives
$\kappa\equiv\kappa_\infty$.

The second equation in
Eq.~\eqref{eq:physical-neutral-orbit-equations} now reduces to
$\partial_a(r\partial_a\varphi)=0$. Multiplication by
$\varphi-\varphi_\infty$ and the identical exhaustion argument give
\begin{equation}
 \int_D r|\nabla_{rz}\varphi|^2\,dr\,dz=0.
 \label{eq:physical-dilaton-energy-zero}
\end{equation}
Thus $\varphi\equiv\varphi_\infty$, proving
Eq.~\eqref{eq:physical-axidilaton-rigidity-conclusion}.
\end{proof}

\begin{remark}[Sharp logical content of the rigidity theorem]
\label{rem:physical-axidilaton-rigidity-scope}
The proof fails, for identifiable reasons, in the presence of gauge flux,
direct matter--scalar coupling, a scalar-charged shell, a horizon or another
inner boundary, a second asymptotic end, axionic monodromy, singularities, or
insufficient decay. These are possible mechanisms for nonconstant scalars;
the theorem does not assert that any of them produces an acceptable galaxy.
\end{remark}

\subsection{Exact Einstein--string-frame dust dictionary}

In four dimensions and in the normalization of
Eq.~\eqref{eq:physical-EMDA-action}, fix the conformal relation
\begin{equation}
 g_E=e^{-2\varphi}g_s,
 \qquad
 \sqrt{-g_E}=e^{-4\varphi}\sqrt{-g_s}.
 \label{eq:physical-frame-relation}
\end{equation}
This is a field redefinition. Which frame is minimally coupled is an
additional physical choice about the matter sector.

\begin{proposition}[Conformal transformation of the Brown dust action]
\label{prop:physical-frame-dust}
Keep the vector density $\mathcal J^\mu$ and the Clebsch variables fixed
under Eq.~\eqref{eq:physical-frame-relation}. Then
\begin{equation}
 |\mathcal J|_E=e^{-\varphi}|\mathcal J|_s,
 \qquad
 n_E=e^{3\varphi}n_s,
 \qquad
 u_E^\mu=e^{\varphi}u_s^\mu,
 \label{eq:physical-frame-current-dictionary}
\end{equation}
and the Einstein-minimal action
Eq.~\eqref{eq:physical-Brown-action} becomes exactly
\begin{equation}
 \mathcal S_{\rm B}^{E}
 =\int_{\mathcal M}\left[
 -m e^{-\varphi}|\mathcal J|_s
 +\mathcal J^\mu
 \left(\partial_\mu\vartheta
 +\beta_A\partial_\mu\alpha^A\right)
 \right]\,d^4x.
 \label{eq:physical-string-frame-dust-action}
\end{equation}
Thus the same particles have string-frame mass
$m_s(\varphi)=me^{-\varphi}$; they are not minimally coupled to $g_s$.

Let
\begin{equation}
 T^{s}_{\mu\nu}:=
 -\frac{2}{\sqrt{-g_s}}
 \frac{\delta\mathcal S_{\rm B}^{E}}{\delta g_s^{\mu\nu}},
 \qquad
 \rho_s:=m e^{-\varphi}n_s.
 \label{eq:physical-string-stress-definition}
\end{equation}
On the dust equations,
\begin{align}
 T_s^{\mu\nu}&=\rho_s u_s^\mu u_s^\nu,
 \label{eq:physical-string-dust-stress}\\
 \nabla^s_\mu T_s^{\mu}{}_{\nu}
 &=\rho_s\partial_\nu\varphi,
 \label{eq:physical-string-stress-nonconservation}\\
 u_s^\mu\nabla^s_\mu u_s^\nu
 &=\left(g_s^{\nu\lambda}+u_s^\nu u_s^\lambda\right)
 \partial_\lambda\varphi.
 \label{eq:physical-string-fifth-force}
\end{align}
Conversely, if constant-mass dust is postulated to be minimal in the string
frame, its Einstein-frame mass is $m_E(\varphi)=me^{\varphi}$, and the
right-hand sides of
Eqs.~\eqref{eq:physical-string-stress-nonconservation} and
~\eqref{eq:physical-string-fifth-force} acquire the opposite sign after
interchanging $s$ and $E$.
\end{proposition}

\begin{proof}
Equation~\eqref{eq:physical-frame-current-dictionary} follows directly
from the definition
$|\mathcal J|^2=-g_{\mu\nu}\mathcal J^\mu\mathcal J^\nu$ and the volume
scaling in Eq.~\eqref{eq:physical-frame-relation}. The current--potential
term is a coordinate density and is conformally unchanged, which proves
Eq.~\eqref{eq:physical-string-frame-dust-action}.

The scalar source of this action is
\begin{equation}
 \frac{1}{\sqrt{-g_s}}
 \frac{\delta\mathcal S_{\rm B}^{E}}{\delta\varphi}=\rho_s.
 \label{eq:physical-string-scalar-source}
\end{equation}
Diffeomorphism invariance therefore gives the Ward identity
Eq.~\eqref{eq:physical-string-stress-nonconservation}. Particle-number
conservation remains $\nabla^s_\mu(n_s u_s^\mu)=0$. Projecting
Eq.~\eqref{eq:physical-string-stress-nonconservation} orthogonally to
$u_s$ proves Eq.~\eqref{eq:physical-string-fifth-force}. Starting instead
from $-m|\mathcal J|_s$ gives
$-me^{\varphi}|\mathcal J|_E$, which proves the converse statement.
\end{proof}

The fifth force in Eq.~\eqref{eq:physical-string-fifth-force} is not a
coordinate artifact: it records which metric defines constant particle
masses and freely falling matter. A conformal rewriting neither creates a
new solution nor preserves the geodesic-dust interpretation in both frames.
In particular, the neutral no-hair theorem above assumes Einstein-minimal
dust; changing the matter coupling changes its scalar equation and its
hypotheses.

\subsection{Conditional hidden symmetries, Buscher duality, and junctions}

Source-free stationary EMDA admits a finite-dimensional hidden
$Sp(4,\mathbb R)$ symmetry, denoted $Sp(2,\mathbb R)$ in part of the physics
literature, and the two-Killing reduction admits a Lax representation and
an affine extension
~\cite{GaltsovKechkin1994PRD,Galtsov1995PRL}. These statements concern the
source-free reduced equations. With dust present, the matrix equation has
the schematic form
\begin{equation}
 \partial_r\!\left(r\mathcal M_r\mathcal M^{-1}\right)
 +\partial_z\!\left(r\mathcal M_z\mathcal M^{-1}\right)
 =\mathcal J_{\rm d}.
 \label{eq:physical-sourced-EMDA-matrix-equation}
\end{equation}
No zero-curvature representation has been established for a generic
nonzero source $\mathcal J_{\rm d}$. A hidden-symmetry transformation may
therefore be applied to a valid vacuum exterior seed, but not presumed to
map the rotating-dust interior to another solution with the same matter.

Buscher duality is a different operation: it acts locally on a string-frame
sigma-model background $(g_s,B,\Phi_s)$ with an Abelian isometry
~\cite{Buscher1987PLB,Buscher1988PLB}. In units in which $\alpha'=1$ and the
sigma-model metric components are dimensionless, let $y$ be the isometry
coordinate and let $i,j$ denote the other coordinates. The Abelian Buscher
rules are
\begin{align}
 \widetilde g_{yy}&=\frac{1}{g_{yy}},
 &
 \widetilde g_{yi}&=\frac{B_{yi}}{g_{yy}},
 &
 \widetilde B_{yi}&=\frac{g_{yi}}{g_{yy}},
 \label{eq:physical-Buscher-y}\\
 \widetilde g_{ij}
 &=g_{ij}-\frac{g_{yi}g_{yj}-B_{yi}B_{yj}}{g_{yy}},
 &
 \widetilde B_{ij}
 &=B_{ij}-\frac{g_{yi}B_{yj}-B_{yi}g_{yj}}{g_{yy}},
 \label{eq:physical-Buscher-transverse}\\
 \widetilde\Phi_s
 &=\Phi_s-\frac12\log g_{yy}.
 \label{eq:physical-Buscher-dilaton}
\end{align}
For the spacelike duality used below they require $g_{yy}>0$ and a
well-defined local isometry quotient.
Restoring dimensions replaces the logarithm by the corresponding
dimensionless ratio in string units.

For illustration only, suppose that the helical plateau
Eq.~\eqref{eq:helical-plateau-metric} is used as a string-frame seed, with
constant dilaton, vanishing initial $B$-field, and the constant
Einstein--string conformal factor absorbed into the normalization of the
displayed metric. Taking $y=\phi$ then gives
$g_{\phi\phi}=r^2-q_0^2$ and $g_{t\phi}=q_0$, and hence
\begin{equation}
 \widetilde B_{\phi t}=\frac{q_0}{r^2-q_0^2},
 \qquad
 \widetilde g_{\phi\phi}=\frac{1}{r^2-q_0^2}.
 \label{eq:physical-Buscher-plateau}
\end{equation}
This conditional dual is singular at the chronology cylinder. More generally, axial
duality at a regular fixed axis has $g_{\phi\phi}\sim r^2$, so
$\widetilde g_{\phi\phi}\sim r^{-2}$ and
$\widetilde\Phi_s\sim\Phi_s-\log r$. It is therefore neither a regular
axis nor a global galactic completion.

The field enlargement also enlarges the junction problem. In the
conventions of Eq.~\eqref{eq:physical-EMDA-action}, introduce the
antisymmetric Maxwell excitation tensor
\begin{equation}
 \mathcal G^{\mu\nu}
 :=e^{-2\varphi}F^{\mu\nu}
 +\kappa{}^\star F^{\mu\nu},
 \qquad
 \nabla_\mu{}^\star F^{\mu\nu}=0,
 \qquad
 \nabla_\mu\mathcal G^{\mu\nu}=0
\label{eq:physical-EMDA-excitation}
\end{equation}
in a source-free neighbourhood of the matching surface.

\begin{proposition}[Necessary shell-free EMDA junction data]
\label{prop:physical-EMDA-junction-data}
Let a non-null hypersurface $\Sigma$ separate two piecewise $C^2$ EMDA
solutions. Assume there is no gravitational shell, scalar surface source,
electric surface current, or magnetic surface current. In a common normal
orientation, the distributional field equations require
\begin{align}
 [h_{ab}]&=0,
 &
 [K_{ab}]&=0,
 \label{eq:physical-EMDA-gravity-junction}\\
 [\varphi]&=0,
 &
 [n(\varphi)]&=0,
 \label{eq:physical-EMDA-dilaton-junction}\\
 [\kappa]&=0,
 &
 [e^{4\varphi}n(\kappa)]&=0,
 \label{eq:physical-EMDA-axion-junction}\\
 [i^*F]&=0,
 &
 [n_\mu\mathcal G^{\mu a}]&=0.
 \label{eq:physical-EMDA-Maxwell-junction}
\end{align}
Here $i:\Sigma\hookrightarrow\mathcal M$ is the inclusion. These are
necessary local data; they do not prove that a global exterior realizing
them exists.
\end{proposition}

\begin{proof}
The first line is Proposition~\ref{prop:physical-Darmois-momentum}. A jump
of either scalar itself would produce an inadmissible derivative of a delta
distribution in its second-order equation, while integration of the scalar
equations across a Gaussian pillbox gives the two normal-flux conditions.
Writing a piecewise smooth two-form as
$F=F^+\mathbf 1_++F^-\mathbf 1_-$ shows that the singular part of $dF$ is
the normal one-form wedged with $[F]$; its vanishing is equivalent to
$[i^*F]=0$. Integrating
$\nabla_\mu\mathcal G^{\mu\nu}=0$ across the same pillbox proves the last
condition. This is the non-null, source-free specialization of the general
distributional scalar-junction framework
~\cite{BarrabesBressange1997CQG}.
\end{proof}

Applying a duality only to the exterior generally violates one or more of
Eqs.~\eqref{eq:physical-EMDA-gravity-junction}--\eqref{eq:physical-EMDA-Maxwell-junction}
and creates a shell or surface charge. Even when the local junction data
pass, asymptotic normalization, NUT charge, gauge charge, scalar charge,
axis regularity, conicity, topology, and causality must all be rechecked.

\begin{remark}[Non-implications]
\label{rem:physical-string-nonimplications}
Nothing in this section implies that the dust interior is an EMDA solution
with nonconstant scalars, that an EMDA hidden symmetry survives a dust
source, or that Buscher duality preserves pressureless geodesic dust. A
pressureless dust element is not a fundamental-string worldsheet or a
string gas; those sources have additional tension data. The helical flat
quotient is not evidence for a string source. A transformed finite-parameter
black-hole exterior is not a generic galactic exterior. Finally, calling
the retained $U(1)$ field electromagnetic or heterotic requires a specified
compactification and matter-coupling prescription. The only unconditional
new result in the string-effective direction is the neutral rigidity
theorem under its stated hypotheses.
\end{remark}
